% Đề mục lấy từ thesis-fixx.tex tại commit c70143bb4e2d3ebef8ab79aab424e9b668bcfe55.
% Biên dịch bằng XeLaTeX.
\documentclass[11pt,a4paper]{article}
\usepackage[top=2.2cm,bottom=2.1cm,left=2.65cm,right=2.45cm]{geometry}
\usepackage{fontspec}
\IfFontExistsTF{texgyretermes-regular.otf}
  {\setmainfont{texgyretermes}[Extension=.otf,UprightFont=*-regular,
    BoldFont=*-bold,ItalicFont=*-italic,BoldItalicFont=*-bolditalic]}
  {\setmainfont{Latin Modern Roman}}
\usepackage[hidelinks]{hyperref}
\hypersetup{pdftitle={Đề cương khóa luận — Nguyễn Trung Nghĩa},
  pdfauthor={Nguyễn Trung Nghĩa}}
\setlength{\parindent}{0pt}
\setlength{\parskip}{0pt}
\raggedbottom
\newcommand{\majorheading}[1]{%
  \par\addvspace{8pt}{\bfseries #1}\par\nopagebreak[4]}
\newcommand{\majorentry}[1]{\par\addvspace{6pt}{\bfseries #1}\par}
\newcommand{\entry}[3]{%
  \par\begingroup\leftskip=#1\relax\hangindent=3.6em
  \makebox[3.6em][l]{#2}#3\par\endgroup}
\newcommand{\methodheading}[1]{\par\hspace*{5em}\texttt{#1}\par}
\begin{document}
\begin{center}
  {\LARGE\bfseries ĐỀ CƯƠNG KHÓA LUẬN\par}
  \vspace{5pt}
  {\large Nghiên cứu mở rộng công cụ USE hỗ trợ đặc tả\par
   và kiểm chứng hệ thống đa tác tử\par}
  \vspace{5pt}
  Nguyễn Trung Nghĩa
\end{center}
\majorentry{MỞ ĐẦU}
\majorheading{CHƯƠNG 1. KIẾN THỨC NỀN TẢNG}
\entry{0pt}{1.1}{Giới thiệu}
\entry{0pt}{1.2}{Ngôn ngữ hướng miền}
\entry{1.35em}{1.2.1}{Khái niệm, phạm vi và lựa chọn DSL}
\entry{1.35em}{1.2.2}{Cú pháp, ngữ nghĩa và metamodel}
\entry{1.35em}{1.2.3}{ANTLR và quy trình xử lý ngôn ngữ}
\entry{0pt}{1.3}{Hệ thống đa tác tử và JaCaMo}
\entry{1.35em}{1.3.1}{Tổng quan MAS và JaCaMo}
\entry{1.35em}{1.3.2}{Mô hình tác tử Jason}
\entry{1.35em}{1.3.3}{Artifact và Workspace trong CArtAgO}
\entry{1.35em}{1.3.4}{Mô hình tổ chức Moise}
\entry{0pt}{1.4}{Mô hình mục tiêu iStar}
\entry{1.35em}{1.4.1}{Actor và các phần tử ý định}
\entry{1.35em}{1.4.2}{Các quan hệ ý định}
\entry{1.35em}{1.4.3}{Góc nhìn và phân tích mô hình iStar}
\entry{0pt}{1.5}{UML, OCL và công cụ USE}
\entry{1.35em}{1.5.1}{Mô hình lớp, đối tượng và trạng thái}
\entry{1.35em}{1.5.2}{OCL và các loại ràng buộc}
\entry{1.35em}{1.5.3}{Môi trường USE và khả năng mở rộng}
\entry{0pt}{1.6}{Kiểm chứng hệ thống đa tác tử}
\entry{1.35em}{1.6.1}{Kiểm chứng mô hình}
\entry{1.35em}{1.6.2}{Kiểm chứng khi thực thi}
\entry{1.35em}{1.6.3}{Trace, snapshot và checkpoint}
\entry{0pt}{1.7}{Các nghiên cứu liên quan}
\entry{0pt}{1.8}{Tổng kết chương}
\newpage
\majorheading{CHƯƠNG 2. PHƯƠNG PHÁP ĐẶC TẢ VÀ KIỂM CHỨNG HỆ THỐNG ĐA TÁC TỬ}
\entry{0pt}{2.1}{Giới thiệu}
\entry{0pt}{2.2}{Động cơ nghiên cứu}
\entry{0pt}{2.3}{Tổng quan phương pháp}
\entry{0pt}{2.4}{Đặc tả CSL trong USE}
\entry{0pt}{2.5}{Đặc tả iStar mở rộng trong USE}
\entry{0pt}{2.6}{Ánh xạ trạng thái JaCaMo sang CSL}
\entry{1.35em}{2.6.1}{Khai báo tương ứng bằng XML}
\entry{1.35em}{2.6.2}{Đồng bộ JaCaMo với CSL}
\entry{1.35em}{2.6.3}{Chiếu đối tượng, thuộc tính và quan hệ}
\entry{0pt}{2.7}{Kiểm chứng ràng buộc và đánh giá Goal/Task}
\entry{1.35em}{2.7.1}{Đánh giá OCL trên trạng thái CSL}
\entry{1.35em}{2.7.2}{Đánh giá Task}
\entry{1.35em}{2.7.3}{Đánh giá Goal}
\entry{1.35em}{2.7.4}{Tổng hợp quyết định checkpoint}
\entry{0pt}{2.8}{Tổng kết chương}
\majorheading{CHƯƠNG 3. THỰC NGHIỆM}
\entry{0pt}{3.1}{Kiến trúc phần mềm}
\entry{1.35em}{3.1.1}{USE trước khi mở rộng}
\entry{1.35em}{3.1.2}{Kiến trúc USE---JaCaMo}
\entry{0pt}{3.2}{Hướng dẫn cài đặt và kịch bản kiểm thử}
\entry{0pt}{3.3}{Tiêu chí kiểm thử}
\entry{0pt}{3.4}{Kết quả đánh giá}
\entry{1.35em}{3.4.1}{Bảo toàn trạng thái}
\entry{1.35em}{3.4.2}{Phản hồi kiểm chứng}
\entry{1.35em}{3.4.3}{Tái hiện kết quả}
\entry{0pt}{3.5}{Thảo luận và hạn chế}
\entry{1.35em}{3.5.1}{Hạn chế của phương pháp}
\entry{1.35em}{3.5.2}{Hướng cải tiến}
\entry{0pt}{3.6}{Kết luận chương}
\majorentry{KẾT LUẬN}
\majorentry{TÀI LIỆU THAM KHẢO}
\majorheading{PHỤ LỤC. MÃ NGUỒN THỰC THI THUẬT TOÁN}
\entry{0pt}{A.1}{Chọn ranh giới checkpoint}
\methodheading{initial}
\methodheading{group}
\methodheading{onMarker}
\methodheading{finish}
\methodheading{send}
\methodheading{onlyReturnsRemain}
\methodheading{goalFailed}
\entry{0pt}{A.2}{Chiếu ảnh chụp nguồn sang trạng thái CSL}
\methodheading{prepare}
\methodheading{applySnapshot}
\methodheading{update}
\end{document}
